\documentclass[11pt,twoside,a4paper]{book}
\usepackage[shorthands=off,english]{babel} % package for multilingual support

\RequirePackage{iftex}
\ifPDFTeX
    \usepackage[utf8]{inputenc} 
    \usepackage[T1]{fontenc}
    % \usepackage{lmodern}
\else
    \RequirePackage{fontspec} % UFT8 fonts for LuaLaTeX
    % \setmainfont{Latin Modern Roman}
\fi
\usepackage{csquotes}

\usepackage[ backend=biber
           , style=numeric
           , sortlocale=en_US
           , bibencoding=UTF8
           , maxcitenames=3
           , maxbibnames=100
           ]{biblatex}

\usepackage{xcolor}
\definecolor{dark-red}{rgb}{0.6,0.15,0.15}
\definecolor{dark-green}{rgb}{0.15,0.4,0.15}
\definecolor{medium-blue}{rgb}{0,0,0.5}
\usepackage[ plainpages=false
           , pdfpagelabels
           , unicode
           , draft=false
           , colorlinks=true
           , linkcolor={dark-red}
           , citecolor={dark-green}
           , urlcolor={medium-blue}
           , unicode=true
           ]{hyperref}

\usepackage{amssymb,amsmath}
\usepackage{verbatim}
\usepackage[final]{listings}
\lstset{ language=C++ }
\usepackage{paralist}
% use upquote for straight quotes in verbatim environments
\usepackage{upquote}
\usepackage{markdownsyntax}

\usepackage{xspace}
\usepackage{pgf}
\usepackage{tikz}
\usetikzlibrary{arrows,positioning}

\usepackage{multirow}
\usepackage{microtype}
\usepackage{tabularx}
\newcolumntype{C}{>{\centering\arraybackslash}X}

\newcommand{\divine}{\mbox{\textsf{DIVINE}}\xspace}
\newcommand{\lart}{\mbox{\textsf{LART}}\xspace}

\newcommand{\paradise}{ParaDiSe\xspace}
\newcommand{\buchi}{büchi\xspace}
\newcommand{\mc}{model checking\xspace}
\newcommand{\mcer}{model-checker\xspace}
\newcommand{\Mcer}{Model-checker\xspace}
\newcommand{\Mc}{Model checking\xspace}

\newcommand{\spin}{\mbox{\textsc{Spin}}\xspace} 
\newcommand{\ltsmin}{\mbox{LTSmin}\xspace}

\newcommand{\llvm}{\textsf{LLVM}\xspace}
\newcommand{\ltl}{LTL\xspace}

\newcommand{\dfs}{DFS\xspace}
\newcommand{\bfs}{BFS\xspace}

\newcommand{\tc}{tree compression\xspace}
\newcommand{\Tc}{Tree compression\xspace}
\newcommand{\tced}{tree-compressed\xspace}
\newcommand{\Tced}{Tree-compressed\xspace}

\newcommand{\darcs}{Darcs\xspace}

\newcommand{\tso}{TSO\xspace}
\newcommand{\pso}{PSO\xspace}
\newcommand{\xtso}{x86-TSO}

\newcommand{\FI}{Faculty of Informatics}
\newcommand{\MU}{Masaryk University}

\newcommand{\Jirik}{doc. RNDr. Jiří Barnat, Ph.D.}
\newcommand{\Mornfall}{RNDr. Petr Ročkai, Ph.D.}

\newcommand{\thesistitle}{\llvm{} Transformations for Model Checking} % enter thesis title
\newcommand{\thesissubtitle}{Master's thesis}
\newcommand{\thesisauthor}{Vladimír Štill}
\newcommand{\thesisYearCity}{Brno, 2016}
\newcommand{\thesisadvisor}{\Jirik}

\newcommand{\note}{\textit{Note:}\xspace}

\addbibresource{bibliography.bib}

\usepackage{setspace}

\oddsidemargin .85in
\evensidemargin .65in
\textwidth 12.7cm

\providecommand{\tightlist}{%
  \setlength{\itemsep}{0pt}\setlength{\parskip}{0pt}}

\newcommand{\TODO}[1]{\textbf{\color{red}#1}}

% autoref names
\addto\extrasenglish{
    \renewcommand{\chapterautorefname}{Chapter}
    \renewcommand{\sectionautorefname}{Section}
    \renewcommand{\subsectionautorefname}{Section}
}

\begin{document}

% initial pages from Mornfall + modifications

\frontmatter
\pagestyle{empty}
 
\begin{center}
    {\Large \sc \FI, \MU}
    \vskip4em
    \includegraphics[width = 4cm, height = 4cm] {logo_fi.pdf}
    \vskip4em
    {\begin{spacing}{1}
        \Huge \bf \thesistitle
    \end{spacing}}
    \vskip2em
    {\Large \sc \thesissubtitle}
    \vskip4em
    {\LARGE \bf \thesisauthor}
    \vfill
    {\hfill \large \thesisYearCity}
\end{center}
 
\cleardoublepage

\section*{Declaration} % from Mornfall
Thereby I declare that this thesis is my original work, which I have
created on my own. All sources and literature used in writing the
thesis, as well as any quoted material, are properly cited, including
full reference to its source.

\vfill
\textbf{Advisor:} \thesisadvisor

\cleardoublepage

\section*{Abstract}

This work focuses on application of \llvm transformations as a preprocessing
step for verification of real-world C and C++ programs using explicit-state
model checker \divine \cite{DiVinE30}. We demonstrate that \llvm transformations
can be used for extension of verifier capabilities and for reduction of the
state space size.

In the case of extension of verifier capabitilies, the main focus is on
verification with relaxed memory models, this is a continuation of the work
started in \cite{SRB15}. We extend the previous transformation to enable
verification of wider range of safety properties, support code with atomic
instructions, support more relaxed memory models than total store order, and
improve state space size of the programs which use this transformation. The
final implementation of this transformation is evaluated and compared with the
previous implementation.

As for state space reductions, we propose the concept of optimizations which
preserve verified property, evaluate some of these optimizations, and propose
additional transformations which can be implemented as future work.

\section*{Keywords}
Formal Verification, C, C++, \llvm, Model Checking, Parallel,
\divine, \lart, \llvm Transformation, Weak Memory Models, Total Store Order,
Implementation.

\cleardoublepage

\section*{Acknowledgements}
I would like to thank all the people in \paradise laboratory for the support and
for the nice work environment. Namely, I would like to thank Jiří Barnat for
advising this thesis and Petr Ročkai for consulting with me.

I would also like to thank my friends and my family for supporting me and having
patience with me.

\cleardoublepage
% \clearpage
\thispagestyle{empty}

\pagestyle{headings}
\tableofcontents % prints table of contents
\mainmatter


\chapter{Introduction}\label{chap:introduction}
\input{intro}

\chapter{\llvm} \label{chap:llvm}
\input{llvm}

\chapter{\divine} \label{chap:divine}
\input{divine}

\chapter{Related Work} \label{chap:related}
\input{related}

\chapter{Proposed \llvm Transformations} \label{chap:transform}
\input{transform}

\chapter{Results} \label{chap:results}
\input{results}

\chapter{Conclusion} \label{chap:conclusion}
\input{conclusion}

\appendix
\chapter{Archive Structure and Compilation} \label{chap:appendix}
\input{appendix}

\chapter*{Bibliography}
\addcontentsline{toc}{chapter}{Bibliography}
\printbibliography[heading=none]

\end{document}
